\chapter*{执行摘要}
\addcontentsline{toc}{chapter}{执行摘要}

智舆面向地级市、县域、园区、景区、高校与中小企业的舆情治理需求，关注传统系统在数据覆盖、人工研判、处置建议与下沉部署成本上的结构性短板。项目提出以多源采集、语义分析、态势研判、3D 可视化与多 Agent 决策推演为主线的 AI 城市治理辅助决策系统，覆盖旅游投诉、民生热线、环境问题、消费维权、校园舆情与企业品牌风险等高频业务场景。系统目标不是扩大数据体量，而是把分散信息转化为可执行的治理洞察。

技术实现依托 AICity 工程基础，形成数据采集、治理知识、模型分析、城市可视化与决策推演五层架构。多模态采集、风险识别、事件抽取、知识聚合、趋势预测与空间表达构成协同闭环，并通过城市模型矩阵与 LoRA 适配器强化区域化知识。相比规则驱动型方案，系统更强调多源融合、空间感知、策略预演与可解释的治理流程，定位为研判人员的辅助工具，而非替代判断。

商业模式采用 SaaS 订阅、轻量私有化、项目制交付与增值服务并行结构，面向不同客户提供基础版、专业版、政务版与定制版。市场切入策略遵循“小切口、强场景、先试点、后复制”，第一年以信阳及周边基层治理与文旅服务为样板，第二年扩展省内地市客户，第三年再向相似场景复制。项目主体为和鸣科技（武汉）有限公司，当前与郑州浩龙机电设备有限公司达成初步合作共识，拟以本地部署方式推进首轮试点，合同截至成稿尚未签署。

可落地性强调真实性优先。当前已具备界面截图、原型展示、功能演示、代码仓库、测试记录、用户沟通材料、试点方案与阶段里程碑等证据；对合同、订单与正式合作材料，统一标注为待补充或待核实，不预设已落地结论。

财务规划以保守测算为前提，通过标准化产品与项目服务结合，逐步建立可复制的区域经营模式。项目价值不仅体现在商业收益，也体现在基层治理效率提升、公共服务响应改善、文旅形象管理优化与学生复合能力培养等社会效益。智舆希望以可信、克制、可迭代的路径形成下沉市场的务实样板。

\begin{infobox}{真实性提醒位}
	\t本书中涉及合同金额、用户数量、样本规模、经营流水等内容，如当前缺乏正式证明材料，均应在附录中标注待补证据类型与补证计划，不得写成已落地既成事实。
\end{infobox}
